\section{Proof of Theorem~\ref{spherical_band}}\label{proof_of_band_comparison}

\subsection[Proof of Theorem~\ref{spherical_band} for n even]{Proof of Theorem~\ref{spherical_band} for $\boldsymbol{n}$ even}

We first prove Theorem~\ref{spherical_band} for even $n$, which is not treated in \cite{Cecchini-Zeidler}.
The necessary changes in the odd-dimensional case will be explained in the next section.

Fix an even integer $n>2$ and a warping function $\rho\colon [\theta_- , \theta_+] \to \R$ such that $(\log \rho)''<0$.
Let~$g_0$ denote the warped product metric in \eqref{g_0}.
Let $M$ be a compact, connected spin manifold of dimension $n$ with boundary $\partial M$.
Let $g$ be a Riemannian metric on $M$.
Suppose that~${\Phi\colon (M,g) \to \big(S^{n-1} \times [\theta_-,\theta_+],g_0\big)}$ is a smooth map satisfying the assumptions of Theorem~\ref{spherical_band}.

Let $\varphi\colon M \to S^{n-1}$ denote the projection of $\Phi$ to the first factor, and let $\Theta\colon M \to [\theta_-,\theta_+]$ denote the projection of $\Phi$ to the second factor. Since $\Phi = (\varphi,\Theta)$ is $1$-Lipschitz, we obtain
\begin{equation} \label{ineq_for_g}
 g \geq {\rm d}\Theta \otimes {\rm d}\Theta + \rho^2(\Theta) \varphi^* g_{S^{n-1}}.
\end{equation}

\begin{Lemma}
\label{gradient_of_Theta}
We have $|\nabla \Theta| \leq 1$, and the inequality is strict unless ${\rm d}\varphi(\nabla \Theta) = 0$.
\end{Lemma}

\begin{proof}
Evaluating the inequality \eqref{ineq_for_g} at the vector $\nabla \Theta$ gives
\[|\nabla \Theta|^2 \geq |\nabla \Theta|^4 + \rho^2(\Theta) |{\rm d}\varphi(\nabla \Theta)|_{g_{S^{n-1}}}^2.\]
From this, the assertion follows easily.
\end{proof}

We write $\partial M = \partial_+ M \cup \partial_- M$, where
\[
\partial_+ M := \partial M \cap \Theta^{-1}(\{\theta_+\}), \qquad \partial_- M := \partial M \cap \Theta^{-1}(\{\theta_-\}).
\]

\begin{Lemma} \label{deg_of_restriction}
Let $\varphi|_{\partial_- M}\colon \partial_- M \to S^{n-1}$ denote the restriction of $\varphi$ to $\partial_- M$.
Then $\deg(\Phi) = \pm \deg(\varphi|_{\partial_- M})$.
\end{Lemma}

\begin{proof}
We can find a smooth function $\hat{\Theta}\colon M \to [\theta_-,\theta_+]$ such that $\hat{\Theta}^{-1}(\{\theta_+\}) = \partial_+ M$, $\hat{\Theta}^{-1}(\{\theta_-\}) = \partial_- M$, and ${\rm d}\hat{\Theta} \neq 0$ at each point on $\partial_+ M \cup \partial_- M$.
Let us define a map $\hat{\Phi}\colon M \to S^{n-1} \times [\theta_-,\theta_+]$ by $\hat{\Phi} = \big(\varphi,\hat{\Theta}\big)$. Clearly, $\deg\big(\hat{\Phi}\big) = \pm \deg(\varphi|_{\partial_- M})$.
Since $\hat{\Phi}$ is homotopic to $\Phi$ relative to $\partial M$, the assertion follows.
\end{proof}

For even $n$, the boundary $\partial M$ is odd-dimensional which is inconvenient for the index calculations.
As in \cite{Llarull}, we remedy the situation by considering products with circles of large radius and sending the radius to infinity.
Let $r$ be a positive real number. We consider the product~${\tilde{M} = M \times S^1}$ equipped with the product metric $\tilde{g} = g + r^2 g_{S^1}$. We write $\partial \tilde{M} = \partial_+ \tilde{M} \cup \partial_- \tilde{M}$, where
\[\partial_+ \tilde{M} := \partial_+ M \times S^1, \qquad \partial_- \tilde{M} := \partial_- M \times S^1.\]

\begin{Lemma} \label{map_to_Sn}
There exists a smooth map
\begin{equation*}
 h\colon\ S^{n-1} \times S^1 \to S^n
\end{equation*}
of degree $\pm 1$ with the property that $h^* g_{S^n} \leq g_{S^{n-1}} + 4 g_{S^1}$.
\end{Lemma}

\begin{proof}
Fix a smooth $2$-Lipschitz function $\beta\colon [-\pi,\pi] \to [-\pi,\pi]$ such that $\beta(t) = -\pi$ for $t \in \bigl[-\pi,-\frac{7\pi}{8}\bigr]$, $\beta(t) = 0$ for $t \in \bigl[-\frac{\pi}{8},\frac{\pi}{8}\bigr]$, and $\beta(t) = \pi$ for $t \in \big[\frac{7\pi}{8},\pi\big]$. Moreover, let us fix a point $a \in S^{n-1}$. We consider the map
\[
 S^{n-1} \times [-\pi,\pi] \to S^n,\qquad (x,t) \mapsto \begin{cases} (\sin \beta(t) x,\cos \beta(t)) & \text{\rm for $t \in [-\pi,0]$}, \\
 (\sin \beta(t) a,\cos \beta(t)) & \text{\rm for $t \in [0,\pi]$.} \end{cases}
 \]
This gives a map $h\colon S^{n-1} \times S^1 \to S^n$ of degree $\pm 1$. Moreover, $h^* g_{S^n} = \sin^2 \beta(t) g_{S^{n-1}} + \beta'(t)^2 g_{S^1}$ for $t \in [-\pi,0]$ and $h^* g_{S^n} = \beta'(t)^2 g_{S^1}$ for $t \in [0,\pi]$.
\end{proof}

In the following, we assume that $h\colon S^{n-1} \times S^1 \to S^n$ is chosen as in Lemma~\ref{map_to_Sn}. We define a smooth map $\tilde{f}\colon \tilde{M} = M \times S^1\to S^n$, $
\tilde{f}(x,t) = h(\varphi(x),t)$ for $x \in M$ and $t \in S^1$.

Choose a spin structure on $M$ and let $S$ denote the spinor bundle over $M$. Furthermore, let $\tilde{S}$ denote the spinor bundle over $\tilde{M} = M \times S^1$, where $S^1$ is equipped with the trivial spin structure $S^1 \times {\rm Spin}(1) \to S^1 \times {\rm SO}(1)$.
Note that with this choice, $\tilde{S}$ is the pull-back of $S$, as a~Clifford-module bundle, under the projection from $\tilde{M} = M \times S^1$ to $M$.

Let $E_0$ denote the spinor bundle of the round sphere $S^n$.
The bundle $E_0$ is equipped with a preferred bundle metric and connection.
Since $n$ is even, we may decompose $E_0$ in the usual way as $E_0 = E_0^+ \oplus E_0^-$, where $E_0^+$ and $E_0^-$ are the $\pm 1$-eigenbundles of the complex volume form.

Next we need an index computation.

\begin{Proposition}[cf.\ Cecchini--Zeidler \cite{Cecchini-Zeidler}]
\label{index_even_dim}
Consider the indices of the following operators:
\begin{itemize}\itemsep=0pt
\item Let $\mathrm{ind}_1$ denote the index of the Dirac operator on $\tilde{S} \otimes \tilde{f}^* E_0^+$ with boundary conditions $u = -{\rm i} \nu \cdot u$ on $\partial_+ \tilde{M}$ and $u = {\rm i} \nu \cdot u$ on $\partial_- \tilde{M}$.
\item Let $\mathrm{ind}_2$ denote the index of the Dirac operator on $\tilde{S} \otimes \tilde{f}^* E_0^+$ with boundary conditions $u = {\rm i} \nu \cdot u$ on $\partial_+ \tilde{M}$ and $u = -{\rm i} \nu \cdot u$ on $\partial_- \tilde{M}$.
\item Let $\mathrm{ind}_3$ denote the index of the Dirac operator on $\tilde{S} \otimes \tilde{f}^* E_0^-$ with boundary conditions $u = -{\rm i} \nu \cdot u$ on $\partial_+ \tilde{M}$ and $u = {\rm i} \nu \cdot u$ on $\partial_- \tilde{M}$.
\item Let $\mathrm{ind}_4$ denote the index of the Dirac operator on $\tilde{S} \otimes \tilde{f}^* E_0^-$ with boundary conditions $u = {\rm i} \nu \cdot u$ on $\partial_+ \tilde{M}$ and $u = -{\rm i} \nu \cdot u$ on $\partial_- \tilde{M}$.
\end{itemize}
Then $\mathrm{ind}_1+\mathrm{ind}_2 = 0$, $\mathrm{ind}_3+\mathrm{ind}_4 = 0$, and $\max \{\mathrm{ind}_1,\mathrm{ind}_2, \mathrm{ind}_3, \mathrm{ind}_4\} > 0$.
\end{Proposition}

\begin{proof}
Since the boundary conditions are adjoint to each other, we obtain $\mathrm{ind}_1+\mathrm{ind}_2 = 0$ and~${\mathrm{ind}_3+\mathrm{ind}_4 = 0}$.

It remains to show that $\max \{\mathrm{ind}_1,\mathrm{ind}_2, \mathrm{ind}_3, \mathrm{ind}_4\} > 0$.
Suppose that this is false.
Then $\mathrm{ind}_1=\mathrm{ind}_2=\mathrm{ind}_3=\mathrm{ind}_4=0$.
We will apply the holographic index theorem in Appendix~\ref{holographic} and the Atiyah--Singer index theorem to show that the assumption $\mathrm{ind}_1=\mathrm{ind}_3=0$ already leads to a contradiction.

The restriction \smash{$\tilde{S}|_{\partial_- \tilde{M}}$} can be identified with the spinor bundle on $\partial_- \tilde{M}$.
We may write \smash{$\tilde{S}|_{\partial_- \tilde{M}} = S^+ \oplus S^-$}, where $S^+$ and $S^-$ denote the eigenbundles of the volume form on $\partial_- \tilde{M}$.
Equivalently, $S^+$ and $S^-$ can be characterized as the eigenbundles of $i\nu$.
This gives the splitting
\[
 \big(\tilde{S} \otimes \tilde{f}^* E_0^+\big)|_{\partial_- \tilde{M}} = \big(S^+ \otimes \big(\tilde{f}|_{\partial_- \tilde{M}}\big)^* E_0^+\big) \oplus \big(S^- \otimes \big(\tilde{f}|_{\partial_- \tilde{M}}\big)^* E_0^+\big).
\]
Similarly,
\[
 \big(\tilde{S} \otimes \tilde{f}^* E_0^-\big)|_{\partial_- \tilde{M}} = \big(S^+ \otimes \big(\tilde{f}|_{\partial_- \tilde{M}}\big)^* E_0^-\big) \oplus \big(S^- \otimes \big(\tilde{f}|_{\partial_- \tilde{M}}\big)^* E_0^-\big).
 \]
Since $\mathrm{ind}_1 =0$, Corollary~\ref{cor:Freed} tells us that the boundary Dirac operator which maps sections of \smash{$S^+ \otimes \big(\tilde{f}|_{\partial_- \tilde{M}}\big)^* E_0^+$} to sections of \smash{$S^- \otimes \big(\tilde{f}|_{\partial_- \tilde{M}}\big)^* E_0^+$} has index $0$.
Similarly, since $\mathrm{ind}_3=0$, the boundary Dirac operator which maps sections of \smash{$S^+ \otimes \big(\tilde{f}|_{\partial_- \tilde{M}}\big)^* E_0^-$} to sections of $\smash{S^- \otimes} \allowbreak\smash{\big(\tilde{f}|_{\partial_- \tilde{M}}\big)^* E_0^-}$ has index $0$.

To obtain a contradiction, we compute the index of the boundary Dirac operators using the Atiyah--Singer index theorem.
Denote the total $\hat{\mathsf{A}}$-class of $\partial_- \tilde{M}$ by $\hat{\mathsf{A}}\big(\partial_- \tilde{M}\big)$. The Chern character of the bundle \smash{$\big(\tilde{f}|_{\partial_- \tilde{M}}\big)^* E_0^+$} is given by the pull-back of $\mathsf{ch}\big(E_0^+\big)$ under \smash{$\tilde{f}|_{\partial_- \tilde{M}}$}. In particular, the Chern character of the bundle \smash{$\big(\tilde{f}|_{\partial_- \tilde{M}}\big)^* E_0^+$} only contains terms in the $0$-th and $n$-th cohomology groups. Since the boundary Dirac operator which maps sections of \smash{$S^+ \otimes \big(\tilde{f}|_{\partial_- \tilde{M}}\big)^* E_0^+$} to sections of \smash{$S^- \otimes \big(\tilde{f}|_{\partial_- \tilde{M}}\big)^* E_0^+$} has index $0$, the Atiyah--Singer index theorem gives
\begin{align}
0
&= \big\langle \hat{\mathsf{A}}\big(\partial_- \tilde{M}\big) \cup \mathsf{ch}\big(\big( \tilde{f}|_{\partial_- \tilde{M}}\big)^* E_0^+\big),\big[\partial_- \tilde{M}\big] \big\rangle\nonumber \\
&= \dim E_0^+ \cdot \big\langle \hat{\mathsf{A}}\big(\partial_- \tilde{M}\big),\big[\partial_- \tilde{M}\big] \big\rangle + \big\langle \mathsf{ch}\big(\big( \tilde{f}|_{\partial_- \tilde{M}}\big)^* E_0^+\big), \big[\partial_- \tilde{M}\big] \big\rangle\nonumber \\
&= \dim E_0^+ \cdot \big\langle \hat{\mathsf{A}}\big(\partial_- \tilde{M}\big),\big[\partial_- \tilde{M}\big] \big\rangle + \deg\big(\tilde{f}|_{\partial_- \tilde{M}}\big) \cdot \big\langle \mathsf{ch}\big(E_0^+\big), [S^n] \big\rangle. \label{index.E_0^+}
\end{align}
Working with $E_0^-$ instead of $E_0^+$, we similarly obtain
\begin{align}
0
&= \big\langle \hat{\mathsf{A}}\big(\partial_- \tilde{M}\big) \cup \mathsf{ch}\big(\big( \tilde{f}|_{\partial_- \tilde{M}}\big)^* E_0^-\big),\big[\partial_- \tilde{M}\big] \big\rangle\nonumber \\
&= \dim E_0^- \cdot \big\langle \hat{\mathsf{A}}\big(\partial_- \tilde{M}\big),\big[\partial_- \tilde{M}\big] \big\rangle + \big\langle \mathsf{ch}\big(\big( \tilde{f}|_{\partial_- \tilde{M}}\big)^* E_0^-\big), \big[\partial_- \tilde{M}\big] \big\rangle \nonumber\\
&= \dim E_0^- \cdot \big\langle \hat{\mathsf{A}}\big(\partial_- \tilde{M}\big),\big[\partial_- \tilde{M}\big] \big\rangle + \deg\big(\tilde{f}|_{\partial_- \tilde{M}}\big) \cdot \big\langle \mathsf{ch}(E_0^-), [S^n] \big\rangle.\label{index.E_0^-}
\end{align}
In the next step, we subtract \eqref{index.E_0^-} from \eqref{index.E_0^+}.
Using the fact that $\dim E_0^+ = \dim E_0^-$, we obtain
\[
 0 = \deg\big(\tilde{f}|_{\partial_- \tilde{M}}\big) \cdot \big\langle \mathsf{ch}\big(E_0^+\big) - \mathsf{ch}(E_0^-), [S^n] \big\rangle.
 \]
It follows from \cite[Proposition~11.24, Chapter~III]{Lawson-Michelsohn} that $\big\langle \mathsf{ch}\big(E_0^+\big) - \mathsf{ch}(E_0^-), [S^n] \big\rangle = \pm \chi(S^n) =\pm 2 \neq 0$ since $n$ is even.
Thus $\deg\big(\tilde{f}|_{\partial_- \tilde{M}}\big)=0$.

By assumption, the map $\Phi\colon M \to S^{n-1} \times [\theta_-,\theta_+]$ has non-zero degree.
Hence, it follows from Lemma~\ref{deg_of_restriction} that the map $\varphi|_{\partial_- M}\colon \partial_- M \to S^{n-1}$ has non-zero degree.
Consequently, the map $\varphi|_{\partial_- M} \times \id\colon \partial_- M \times S^1 \to S^{n-1} \times S^1$ has non-zero degree. By Lemma~\ref{map_to_Sn}, the map~${h\colon S^{n-1} \times S^1 \to S^n}$ has non-zero degree.
Since \smash{$\tilde{f}|_{\partial_- \tilde{M}} = h \circ (\varphi|_{\partial_- M} \times \id)$}, we conclude that the map \smash{$\tilde{f}|_{\partial_- \tilde{M}}\colon \partial_- \tilde{M} \to S^n$} has non-zero degree.
This is a contradiction.
\end{proof}

By Proposition~\ref{index_even_dim}, we know that $\max \{\mathrm{ind}_1,\mathrm{ind}_2, \mathrm{ind}_3, \mathrm{ind}_4\} > 0$. After switching the bundles~$E_0^+$ and $E_0^-$ if necessary, we may assume that $\max \{\mathrm{ind}_1,\mathrm{ind}_2\} > 0$.
In the remainder of this section, we focus on the case $\mathrm{ind}_1 > 0$.
(The case $\mathrm{ind}_2 > 0$ can be treated analogously.)

Let $\tilde{E}$ denote the pull-back of $E_0^+$ under the map $\tilde{f}$. The bundle metric on $E_0^+$ gives us a~bundle metric on $\tilde{E}$. Moreover, the connection on $E_0^+$ induces a connection on $\tilde{E}$. We denote by \smash{$\nabla^{\tilde{S} \otimes \tilde{E}}$} the tensor product connection on $\tilde{S} \otimes \tilde{E}$. We denote by \smash{$\mathcal{D}^{\tilde{S} \otimes \tilde{E}}$} the Dirac operator acting on sections of $\tilde{S} \otimes \tilde{E}$,
\[
 \mathcal{D}^{\tilde{S} \otimes \tilde{E}} u = \sum_{k=1}^{n+1} e_k \cdot \nabla_{e_k}^{\tilde{S} \otimes \tilde{E}} u,
\]
where $\{e_1, \ldots, e_{n+1}\}$ is a local orthonormal frame on $\tilde{M}$.
Finally, we define the boundary Dirac operator by
\[
 \mathcal{D}^{\partial \tilde{M}} u = \sum_{k=1}^n \nu \cdot e_k \cdot \nabla_{e_k}^{\tilde{S} \otimes \tilde{E}} u + \frac{1}{2} H u,
\]
where $\{e_1,\hdots,e_n\}$ is a local orthonormal frame on $\partial \tilde{M}$.
The boundary Dirac operator is self-adjoint and anti-commutes with Clifford multiplication by $\nu$.

Recall the Weitzenb\"ock formula (see \cite[Theorem~8.17, Chapter~II]{Lawson-Michelsohn}),
\[
 \big(\mathcal{D}^{\tilde{S} \otimes \tilde{E}}\big)^2 u
 =
 \big(\nabla^{\tilde{S} \otimes \tilde{E}}\big)^*\nabla^{\tilde{S} \otimes \tilde{E}} u + \frac{1}{4} R u + \mathcal{R}^{\tilde{E}} u,
\]
where $\mathcal{R}^{\tilde{E}}$ is a section of the endomorphism bundle of $\tilde{S} \otimes \tilde{E}$ which depends on the curvature of the bundle $\tilde{E}$.

We define a vector field $T$ on $\tilde{M}$ by $T = \frac{1}{r} \frac{\partial}{\partial t}$, where $t \mapsto (\cos t,\sin t)$ is the canonical local coordinate on $S^1$. Note that $T$ is parallel and has unit length with respect to the metric $\tilde{g}$. In the following, $\Psi$ will denote a smooth function on $M$ which will be specified later. We may extend $\Psi$ to a smooth function on $\tilde{M}$ satisfying $T(\Psi) = 0$. If $u$ is a section of the bundle $\tilde{S} \otimes \tilde{E}$ and $X$ is a vector field on $\tilde{M}$, we define
\begin{equation} \label{def:tildeP}
\tilde{P}_X u = \nabla_{X - \langle X,T \rangle T}^{\tilde{S} \otimes \tilde{E}} u + \frac{\rm i}{2} \Psi \cdot (X - \langle X,T \rangle T) \cdot u.
\end{equation}
Our argument is based on the following integral formula which links several geometric quantities on $\tilde{M}$ and $\partial \tilde{M}$.

\begin{Proposition}
\label{integral_formula_even_dim}
Let $u \in C^\infty\big(\tilde{M},\tilde{S} \otimes \tilde{E}\big)$. Then
\begin{align*}
& - \int_{\tilde{M}} \left|\mathcal{D}^{\tilde{S} \otimes \tilde{E}} u - \frac{{\rm i}n}{2} \Psi u\right|^2 + \int_{\tilde{M}} \big|\tilde{P} u\big|^2 + \int_{\tilde{M}} \big|\nabla_T^{\tilde{S} \otimes \tilde{E}} u\big|^2 + \frac{1}{4} \int_{\tilde{M}} R |u|^2 \\
&+ \int_{\tilde{M}} \big\langle \mathcal{R}^{\tilde{E}} u,u \big\rangle + \frac{n(n-1)}{4} \int_{\tilde{M}} \Psi^2 |u|^2 - \frac{{\rm i}(n-1)}{2} \int_{\tilde{M}} \langle (\nabla \Psi) \cdot u,u \rangle \\
&\qquad= \frac{1}{2} \int_{\partial_+ \tilde{M}} \big\langle \mathcal{D}^{\partial \tilde{M}} u,u + {\rm i} \nu \cdot u \big\rangle + \frac{1}{2} \int_{\partial_+ \tilde{M}} \big\langle u + {\rm i} \nu \cdot u,\mathcal{D}^{\partial \tilde{M}} u \big\rangle \\
&\phantom{\qquad=}{} - \frac{1}{2} \int_{\partial_+ \tilde{M}} (H-(n-1)\Psi) |u|^2 - \frac{n-1}{2} \int_{\partial_+ \tilde{M}} \Psi \langle u + {\rm i} \nu \cdot u,u \rangle \\
&\phantom{\qquad=}{} + \frac{1}{2} \int_{\partial_- \tilde{M}} \big\langle \mathcal{D}^{\partial \tilde{M}} u,u - {\rm i} \nu \cdot u \big\rangle + \frac{1}{2} \int_{\partial_- \tilde{M}} \big\langle u - {\rm i} \nu \cdot u,\mathcal{D}^{\partial \tilde{M}} u \big\rangle \\
&\phantom{\qquad=}{} - \frac{1}{2} \int_{\partial_- \tilde{M}} (H+(n-1)\Psi) |u|^2 + \frac{n-1}{2} \int_{\partial_- \tilde{M}} \Psi \langle u - {\rm i} \nu \cdot u,u \rangle.
\end{align*}
\end{Proposition}

\begin{proof}
Integrating the Weitzenb\"ock formula and using the divergence theorem gives
\begin{align*}
&- \int_{\tilde{M}} \big|\mathcal{D}^{\tilde{S} \otimes \tilde{E}} u\big|^2 + \int_{\tilde{M}} \big|\nabla^{\tilde{S} \otimes \tilde{E}} u\big|^2 + \frac{1}{4} \int_{\tilde{M}} R |u|^2 + \int_{\tilde{M}} \big\langle \mathcal{R}^{\tilde{E}} u,u \big\rangle \\
&\qquad= \int_{\partial \tilde{M}} \big\langle \nu \cdot \mathcal{D}^{\tilde{S} \otimes \tilde{E}} u,u \big\rangle + \int_{\partial \tilde{M}} \big\langle \nabla^{\tilde{S} \otimes \tilde{E}}_\nu u,u \big\rangle.
\end{align*}
Note that $\nu \cdot \mathcal{D}^{\tilde{S} \otimes \tilde{E}} u + \nabla^{\tilde{S} \otimes \tilde{E}}_\nu u = \mathcal{D}^{\partial \tilde{M}} u - \frac{1}{2} H u$.
This gives
\begin{align*}
&- \int_{\tilde{M}} \big|\mathcal{D}^{\tilde{S} \otimes \tilde{E}} u\big|^2 + \int_{\tilde{M}} \big|\nabla^{\tilde{S} \otimes \tilde{E}} u\big|^2 + \frac{1}{4} \int_{\tilde{M}} R |u|^2 + \int_{\tilde{M}} \big\langle \mathcal{R}^{\tilde{E}} u,u \big\rangle \\
&\qquad= \int_{\partial \tilde{M}} \big\langle \mathcal{D}^{\partial \tilde{M}} u,u \big\rangle - \frac{1}{2} \int_{\partial \tilde{M}} H |u|^2 \\
&\qquad= \frac{1}{2} \int_{\partial \tilde{M}} \big\langle \mathcal{D}^{\partial \tilde{M}} u,u \big\rangle + \frac{1}{2} \int_{\partial \tilde{M}} \big\langle u,\mathcal{D}^{\partial \tilde{M}} u \big\rangle - \frac{1}{2} \int_{\partial \tilde{M}} H |u|^2.
\end{align*}
Since $\mathcal{D}^{\partial \tilde{M}}$ is self-adjoint and anti-commutes with $\nu$, we find
\begin{align*}
\int_{\partial_\pm \tilde{M}} \big\langle \mathcal{D}^{\partial \tilde{M}} u,{\rm i} \nu \cdot u \big\rangle
&=
\int_{\partial_\pm \tilde{M}} \big\langle u,\mathcal{D}^{\partial \tilde{M}}({\rm i} \nu \cdot u) \big\rangle =
-\int_{\partial_\pm \tilde{M}} \big\langle u,{\rm i} \nu \cdot\mathcal{D}^{\partial \tilde{M}} u \big\rangle \\
&=
-\int_{\partial_\pm \tilde{M}} \big\langle {\rm i} \nu \cdot u,\mathcal{D}^{\partial \tilde{M}} u \big\rangle .
\end{align*}
Therefore,
\begin{align}
& - \int_{\tilde{M}}\big|\mathcal{D}^{\tilde{S} \otimes \tilde{E}} u\big|^2 + \int_{\tilde{M}} \big|\nabla^{\tilde{S} \otimes \tilde{E}} u\big|^2 + \frac{1}{4} \int_{\tilde{M}} R |u|^2 + \int_{\tilde{M}} \big\langle \mathcal{R}^{\tilde{E}} u,u \big\rangle \nonumber\\
&\qquad= \frac{1}{2} \int_{\partial_+ \tilde{M}} \big\langle \mathcal{D}^{\partial \tilde{M}} u,u + {\rm i} \nu \cdot u \big\rangle + \frac{1}{2} \int_{\partial_+ \tilde{M}} \big\langle u + {\rm i} \nu \cdot u,\mathcal{D}^{\partial \tilde{M}} u \big\rangle\nonumber \\
&\qquad\phantom{=}{} + \frac{1}{2} \int_{\partial_- \tilde{M}} \big\langle \mathcal{D}^{\partial \tilde{M}} u,u - {\rm i} \nu \cdot u \big\rangle + \frac{1}{2} \int_{\partial_- \tilde{M}} \big\langle u - {\rm i} \nu \cdot u,\mathcal{D}^{\partial \tilde{M}} u \big\rangle
 - \frac{1}{2} \int_{\partial \tilde{M}} H |u|^2.\label{eq:formula1}
\end{align}
Using the definition of $\tilde{P} u$ and a local orthonormal frame $e_1,\dots,e_n,T$ on $\tilde{M}$, we compute
\begin{align}
\big|\tilde{P} u\big|^2
={}&
\sum_{k=1}^n \left| \nabla_{e_k}^{\tilde{S} \otimes \tilde{E}} u + \frac{\rm i}{2} \Psi e_k \cdot u \right|^2 \nonumber\\
={}&
\big|\nabla^{\tilde{S} \otimes \tilde{E}} u\big|^2 - \big|\nabla_T^{\tilde{S} \otimes \tilde{E}} u\big|^2 + \frac{n}{4} \Psi^2 |u|^2 \nonumber\\
& + \frac{\rm i}{2} \Psi \big\langle \mathcal{D}^{\tilde{S} \otimes \tilde{E}} u - T \cdot \nabla_T^{\tilde{S} \otimes \tilde{E}} u,u \big\rangle - \frac{\rm i}{2} \Psi \big\langle u,\mathcal{D}^{\tilde{S} \otimes \tilde{E}} u - T \cdot \nabla_T^{\tilde{S} \otimes \tilde{E}} u \big\rangle \nonumber\\
={}&
\big|\nabla^{\tilde{S} \otimes \tilde{E}} u\big|^2 - \big|\nabla_T^{\tilde{S} \otimes \tilde{E}} u\big|^2 + \frac{n}{4} \Psi^2 |u|^2 \nonumber\\
& + \frac{\rm i}{2} \Psi \big\langle \mathcal{D}^{\tilde{S} \otimes \tilde{E}} u,u \big\rangle - \frac{\rm i}{2} \Psi \big\langle u,\mathcal{D}^{\tilde{S} \otimes \tilde{E}} u \big\rangle - \frac{\rm i}{2} \Psi T(\langle T \cdot u,u \rangle).\label{eq:Pu2}
\end{align}
Using the divergence theorem, we find
\begin{align}
\int_{\tilde{M}} \Psi T(\langle T \cdot u,u \rangle)
={}& \int_{\partial \tilde{M}} \Psi \langle T \cdot u,u \rangle \langle T,\nu \rangle \nonumber\\
& - \int_{\tilde{M}} \Psi \langle T \cdot u,u \rangle \mathrm{div} T - \int_{\tilde{M}} T(\Psi) \langle T \cdot u,u \rangle = 0 .\label{eq:lastterm}
\end{align}
We integrate \eqref{eq:Pu2} over $\tilde{M}$ and insert \eqref{eq:lastterm} and obtain
\begin{align}
\int_{\tilde{M}} \big|\tilde{P} u\big|^2
={}& \int_{\tilde{M}} \big|\nabla^{\tilde{S} \otimes \tilde{E}} u\big|^2 - \int_{\tilde{M}} \big|\nabla_T^{\tilde{S} \otimes \tilde{E}} u\big|^2 + \frac{n}{4} \int_{\tilde{M}} \Psi^2 |u|^2\nonumber \\
& + \frac{\rm i}{2} \int_{\tilde{M}} \Psi \big\langle \mathcal{D}^{\tilde{S} \otimes \tilde{E}} u,u \big\rangle - \frac{\rm i}{2} \int_{\tilde{M}} \Psi \big\langle u,\mathcal{D}^{\tilde{S} \otimes \tilde{E}} u \big\rangle.\label{eq:Pu2neu}
\end{align}
Substituting \eqref{eq:Pu2neu} into \eqref{eq:formula1}, we obtain
\begin{align}
&- \int_{\tilde{M}} \big|\mathcal{D}^{\tilde{S} \otimes \tilde{E}} u\big|^2 + \int_{\tilde{M}} \big|\tilde{P} u\big|^2 + \int_{\tilde{M}} \big|\nabla_T^{\tilde{S} \otimes \tilde{E}} u\big|^2 + \frac{1}{4} \int_{\tilde{M}} R |u|^2
+ \int_{\tilde{M}} \big\langle \mathcal{R}^{\tilde{E}} u,u \big\rangle \nonumber\\
&- \frac{n}{4} \int_{\tilde{M}} \Psi^2 |u|^2
 - \frac{\rm i}{2} \int_{\tilde{M}} \Psi \big\langle \mathcal{D}^{\tilde{S} \otimes \tilde{E}} u,u \big\rangle + \frac{\rm i}{2} \int_{\tilde{M}} \Psi \big\langle u,\mathcal{D}^{\tilde{S} \otimes \tilde{E}} u \big\rangle\nonumber \\
&\qquad= \frac{1}{2} \int_{\partial_+ \tilde{M}} \big\langle \mathcal{D}^{\partial \tilde{M}} u,u + {\rm i} \nu \cdot u \big\rangle + \frac{1}{2} \int_{\partial_+ \tilde{M}} \big\langle u + {\rm i} \nu \cdot u,\mathcal{D}^{\partial \tilde{M}} u \big\rangle\nonumber \\
&\qquad\phantom{=}{} + \frac{1}{2} \int_{\partial_- \tilde{M}} \big\langle \mathcal{D}^{\partial \tilde{M}} u,u - {\rm i} \nu \cdot u \big\rangle + \frac{1}{2} \int_{\partial_- \tilde{M}} \big\langle u - {\rm i} \nu \cdot u,\mathcal{D}^{\partial \tilde{M}} u \big\rangle - \frac{1}{2} \int_{\partial \tilde{M}} H |u|^2.\label{eq:formula2}
\end{align}
Using the divergence theorem, we obtain
\begin{align}
&- \frac{{\rm i}(n-1)}{2} \int_{\tilde{M}} \Psi \big\langle \mathcal{D}^{\tilde{S} \otimes \tilde{E}} u,u \big\rangle + \frac{{\rm i}(n-1)}{2} \int_{\tilde{M}} \Psi \big\langle u,\mathcal{D}^{\tilde{S} \otimes \tilde{E}} u \big\rangle
- \frac{{\rm i}(n-1)}{2} \int_{\tilde{M}} \langle (\nabla \Psi) \cdot u,u \rangle\nonumber \\
&\qquad= - \frac{{\rm i}(n-1)}{2} \int_{\tilde{M}} \big\langle \mathcal{D}^{\tilde{S} \otimes \tilde{E}} (\Psi u),u \big\rangle + \frac{{\rm i}(n-1)}{2} \int_{\tilde{M}} \big\langle \Psi u,\mathcal{D}^{\tilde{S} \otimes \tilde{E}} u \big\rangle\nonumber \\
&\qquad= -\frac{{\rm i}(n-1)}{2} \int_{\partial \tilde{M}} \langle \nu \cdot (\Psi u),u \rangle.\label{eq:formula3}
\end{align}
Adding \eqref{eq:formula2} and \eqref{eq:formula3} gives
\begin{align*}
&- \int_{\tilde{M}} \big|\mathcal{D}^{\tilde{S} \otimes \tilde{E}} u\big|^2 + \int_{\tilde{M}} \big|\tilde{P} u\big|^2 + \int_{\tilde{M}} \big|\nabla_T^{\tilde{S} \otimes \tilde{E}} u\big|^2 + \frac{1}{4} \int_{\tilde{M}} R |u|^2
+ \int_{\tilde{M}} \big\langle \mathcal{R}^{\tilde{E}} u,u \big\rangle \\
&- \frac{n}{4} \int_{\tilde{M}} \Psi^2 |u|^2 - \frac{{\rm i}(n-1)}{2} \int_{\tilde{M}} \langle (\nabla \Psi) \cdot u,u \rangle
- \frac{{\rm i}n}{2} \int_{\tilde{M}} \Psi \big\langle \mathcal{D}^{\tilde{S} \otimes \tilde{E}} u,u \big\rangle + \frac{{\rm i}n}{2} \int_{\tilde{M}} \Psi \big\langle u,\mathcal{D}^{\tilde{S} \otimes \tilde{E}} u \big\rangle\\
&\qquad= \frac{1}{2} \int_{\partial_+ \tilde{M}} \big\langle \mathcal{D}^{\partial \tilde{M}} u,u + {\rm i} \nu \cdot u \big\rangle + \frac{1}{2} \int_{\partial_+ \tilde{M}} \big\langle u + {\rm i} \nu \cdot u,\mathcal{D}^{\partial \tilde{M}} u \big\rangle \\
&\phantom{\qquad=}{} - \frac{1}{2} \int_{\partial_+ \tilde{M}} (H-(n-1)\Psi) |u|^2 - \frac{n-1}{2} \int_{\partial_+ \tilde{M}} \Psi \langle u + {\rm i} \nu \cdot u,u \rangle \\
&\phantom{\qquad=}{} + \frac{1}{2} \int_{\partial_- \tilde{M}} \big\langle \mathcal{D}^{\partial \tilde{M}} u,u - {\rm i} \nu \cdot u \big\rangle + \frac{1}{2} \int_{\partial_- \tilde{M}} \big\langle u - {\rm i} \nu \cdot u,\mathcal{D}^{\partial \tilde{M}} u \big\rangle \\
&\phantom{\qquad=}{} - \frac{1}{2} \int_{\partial_- \tilde{M}} (H+(n-1)\Psi) |u|^2 + \frac{n-1}{2} \int_{\partial_- \tilde{M}} \Psi \langle u - {\rm i} \nu \cdot u,u \rangle.
\end{align*}
This completes the proof of Proposition~\ref{integral_formula_even_dim}.
\end{proof}

At this point, we specify our choice of the function $\Psi$. We define a function $\psi\colon [\theta_-, \theta_+] \to \R$ by $\psi(\theta) := \frac{\rho'(\theta)}{\rho(\theta)}$.
Since $(\log \rho)''<0$, we know that $-\psi' > 0$ on $[\theta_-,\theta_+]$.
Using \eqref{scal_g_0}, the inequality $R \geq R_{g_0} \circ \Phi$ gives
\begin{equation} \label{ineq_for_scal_even_dim}
R \geq (n-1) \left( -2 \psi'(\Theta) - n \psi^2(\Theta) + \frac{n-2}{\rho^2(\Theta)} \right).
\end{equation}
We define $\Psi\colon M \to \R$ by $\Psi = \psi \circ \Theta$.

\begin{Proposition}
\label{existence_of_spinor_even_dim}
Assume that \smash{$r > 2 \sup_{[\theta_-,\theta_+]} \rho$}. Then we can find an element $t_0 \in S^1$ and a~section $u \in C^\infty\big(\tilde{M},\tilde{S} \otimes \tilde{E}\big)$ such that
\[\int_{M \times \{t_0\}} \frac{1}{\rho(\Theta)} |u|^2 = 1\]
and
\[\int_{M \times \{t_0\}} \big|\tilde{P} u\big|^2 \leq \frac{n-1}{r}.\]
\end{Proposition}

\begin{proof}
Recall that we are assuming $\mathrm{ind}_1 > 0$. In view of the deformation invariance of the index, we can find a section $u \in C^\infty\big(\tilde{M},\tilde{S} \otimes \tilde{E}\big)$ such that
\begin{itemize}\itemsep=0pt
 \item $u$ does not vanish identically,
 \item $\mathcal{D}^{\tilde{S} \otimes \tilde{E}} u - \frac{{\rm i}n}{2} \Psi u = 0$ on $\tilde{M}$,
 \item $u = -{\rm i} \nu \cdot u$ on $\partial_+ \tilde{M}$ and $u = {\rm i} \nu \cdot u$ on $\partial_- \tilde{M}$.
\end{itemize}
Using Proposition~\ref{integral_formula_even_dim}, we obtain
\begin{align}
&\int_{\tilde{M}} \big|\tilde{P} u\big|^2 + \int_{\tilde{M}} \big|\nabla_T^{\tilde{S} \otimes \tilde{E}} u\big|^2 + \frac{1}{4} \int_{\tilde{M}} R |u|^2\nonumber \\
&\qquad{}+ \int_{\tilde{M}} \big\langle \mathcal{R}^{\tilde{E}} u,u \big\rangle + \frac{n(n-1)}{4} \int_{\tilde{M}} \Psi^2 |u|^2 - \frac{{\rm i}(n-1)}{2} \int_{\tilde{M}} \langle (\nabla \Psi) \cdot u,u \rangle \nonumber\\
&\phantom{\qquad+}{}= -\frac{1}{2} \int_{\partial_+ \tilde{M}} (H-(n-1)\Psi) |u|^2 - \frac{1}{2} \int_{\partial_- \tilde{M}} (H+(n-1)\Psi) |u|^2.\label{important_estimate_even_dim}
\end{align}
By assumption and using \eqref{mean_g_0}, \[H-(n-1)\Psi = H-(n-1) \frac{\rho'(\theta_+)}{\rho(\theta_+)} \geq 0\] on $\partial_+ \tilde{M}$ and \[H +(n-1)\Psi \geq H+(n-1) \frac{\rho'(\theta_-)}{\rho(\theta_-)} \geq 0\] on $\partial_- \tilde{M}$. Consequently, the right-hand side in \eqref{important_estimate_even_dim} is non-positive.

We next analyze the term $\mathcal{R}^{\tilde{E}}$. To that end, we fix a point $(x,t) \in \tilde{M}$. Let $ \mu_1,\hdots,\mu_{n+1} \geq 0$ denote the singular values of the differential \smash{$d\tilde{f}_{(x,t)}\colon \big(T_{(x,t)} \tilde{M},\tilde{g}\big) \to (T_{\tilde{f}(x,t)} S^n,g_{S^n})$}, arranged in decreasing order. Since the differential \smash{$d\tilde{f}_{(x,t)}$} has rank at most $n$, it follows that $\mu_{n+1} = 0$.
The eigenvalues of the symmetric bilinear form $\tilde{f}^* g_{S^n}$ with respect to the metric $\tilde{g} = g + r^2 g_{S^1}$ are given by $\mu_1^2,\hdots,\mu_n^2,0$.

Using \eqref{ineq_for_g}, we obtain, on the one hand,
\[
 \tilde{g} = g + r^2 g_{S^1} \geq \rho^2(\Theta) \varphi^* g_{S^{n-1}} + r^2 g_{S^1}.
 \]
On the other hand, the inequality $h^* g_{S^n} \leq g_{S^{n-1}} + 4 g_{S^1}$ implies
\[
 \tilde{f}^* g_{S^n} \leq \varphi^* g_{S^{n-1}} + 4 g_{S^1}.
\]
By assumption, $r > 2\rho(\Theta)$ at the point $(x,t)$. Hence, the min-max characterization of the eigenvalues implies that $\mu_1^2,\hdots,\mu_{n-1}^2 \leq \frac{1}{\rho^2(\Theta)}$ and $\mu_n^2 \leq \frac{4}{r^2}$.
Together with Proposition~\ref{curvature.term}, this implies
\begin{align*}
\big\langle \mathcal{R}^{\tilde{E}} u,u \big\rangle
&\geq -\frac{1}{4} \sum_{\substack{1\leq k,j\leq n\\ j \neq k}} \mu_j \mu_k |u|^2
\geq -\frac{(n-2)(n-1)}{4} \frac{1}{\rho^2(\Theta)} |u|^2 - \frac{n-1}{r} \frac{1}{\rho(\Theta)} |u|^2.
\end{align*}
Recall that $|\nabla \Theta| \leq 1$ by Lemma~\ref{gradient_of_Theta}.
Since $-\psi' > 0$ on $[\theta_-,\theta_+]$, it follows that $|\nabla \Psi| \leq -\psi'(\Theta)$.
Using \eqref{ineq_for_scal_even_dim}, we obtain the pointwise estimate
\begin{align*}
&\frac{1}{4} R |u|^2 + \big\langle \mathcal{R}^{\tilde{E}} u,u \big\rangle + \frac{n(n-1)}{4} \Psi^2 |u|^2 - \frac{n-1}{2} |\nabla \Psi| |u|^2 \\
&\qquad\geq \frac{1}{4} R |u|^2 - \frac{(n-2)(n-1)}{4} \frac{1}{\rho^2(\Theta)} |u|^2 - \frac{n-1}{r} \frac{1}{\rho(\Theta)} |u|^2\\
&\phantom{\qquad\geq }{}+ \frac{n(n-1)}{4} \psi^2(\Theta) |u|^2 + \frac{n-1}{2} \psi'(\Theta) |u|^2 \\
&\qquad\geq -\frac{n-1}{r} \frac{1}{\rho(\Theta)} |u|^2.
\end{align*}
Putting these facts together, we conclude that
\[
\int_{\tilde{M}} \big|\tilde{P} u\big|^2 \leq \frac{n-1}{r} \int_{\tilde{M}} \frac{1}{\rho(\Theta)} |u|^2.
\]
Hence, we can find an element $t_0 \in S^1$ such that
\[
\int_{M \times \{t_0\}} \big|\tilde{P} u\big|^2 \leq \frac{n-1}{r} \int_{M \times \{t_0\}} \frac{1}{\rho(\Theta)} |u|^2
\]
and $\int_{M \times \{t_0\}} \frac{1}{\rho(\Theta)} |u|^2 > 0$.
From this, the assertion follows.
\end{proof}


\begin{Corollary}
\label{limiting_spinor_even_dim}
There exists an element $t_0 \in S^1$ with the following property.
Let $f\colon M \to S^n$ be defined by
\[
 f(x) := \tilde{f}(x,t_0) = h(\varphi(x) ,t_0).
\]
Let $E$ denote the pull-back of $E_0^+$ under the map $f$.
Then there exists a section $s \in C^\infty(M,S \otimes E)$ such that
\[
 \int_M |s|^2 = 1
\]
and
\[
 \nabla_X^{S \otimes E} s + \frac{\rm i}{2} \Psi X \cdot s = 0
\]
for every vector field $X$.
\end{Corollary}

\begin{proof}
Let us consider a sequence $r_\ell \to \infty$. For each $\ell$, Proposition~\ref{existence_of_spinor_even_dim} implies the existence of an element $t_\ell \in S^1$ and a section $u^{(\ell)} \in C^\infty\big(\tilde{M},\tilde{S} \otimes \tilde{E}\big)$ such that
\[
 \int_{M \times \{t_\ell\}} \frac{1}{\rho(\Theta)}\big|u^{(\ell)}\big|^2 = 1
\]
and
\[\int_{M \times \{t_\ell\}} \big|\tilde{P} u^{(\ell)}\big|^2 \leq \frac{n-1}{r_\ell}.
\]
After passing to a subsequence, we may assume that the sequence $t_\ell$ converges to an element $t_0 \in S^1$. We define maps $f\colon M \to S^n$ and $f^{(\ell)}\colon M \to S^n$ by
\[
 f(x) := \tilde{f}(x,t_0) = h(\varphi(x),t_0), \qquad f^{(\ell)}(x) := \tilde{f}(x,t_\ell) = h(\varphi(x) ,t_\ell)
\]
for $x \in M$.
Let $E$ denote the pull-back of $E_0^+$ under $f$, and let $E^{(\ell)}$ denote the pull-back of $E_0^+$ under $f^{(\ell)}$.
The restriction of $u^{(\ell)}$ to $M \times \{t_\ell\}$ gives a section $s^{(\ell)} \in C^\infty\big(M,S \otimes E^{(\ell)}\big)$ such that
\[
 \int_M \frac{1}{\rho(\Theta)}\big|s^{(\ell)}\big|^2 = 1
 \]
and
\begin{equation}
 \int_M \sum_{k=1}^n \left| \nabla_{e_k}^{S \otimes E^{(\ell)}} s^{(\ell)} + \frac{\rm i}{2} \psi(\Theta) e_k \cdot s^{(\ell)} \right|^2 \leq \frac{n-1}{r_\ell}.
\label{eq:almostflat}
\end{equation}
For each $\ell$, we define a bundle map $\sigma^{(\ell)}\colon E^{(\ell)} \to E$ as follows.
For each point $x \in M$, the map \smash{$\sigma^{(\ell)}_x\colon E^{(\ell)}_x = \big(E_0^+\big)_{f^{(\ell)}(x)} \to E_x = \big(E_0^+\big)_{f(x)}$} is defined as the parallel transport along the shortest geodesic from $f^{(\ell)}(x) \in S^n$ to $f(x) \in S^n$.
It is easy to see that $\sigma^{(\ell)}$ is a bundle isometry for each~$\ell$. It follows that the map $\big(\id \otimes \sigma^{(\ell)}\big)\colon S \otimes E^{(\ell)} \to S \otimes E$ is a bundle isometry for each $\ell$.

We may write
\begin{equation}
\label{eq:A_l_term}
\nabla^{S \otimes E}\big(\big(\id \otimes \sigma^{(\ell)}\big) s^{(\ell)}\big) = \big(\id \otimes \sigma^{(\ell)}\big) \big(\nabla^{S \otimes E^{(\ell)}} s^{(\ell)} + A^{(\ell)} s^{(\ell)}\big),
\end{equation}
where $A^{(\ell)}$ is a $1$-form taking values in the endomorphism bundle $\text{\rm End}\big(S \otimes E^{(\ell)}\big)$. Since $t_\ell \to t_0$, the maps $f_\ell$ converge to $f$ smoothly. From this, we deduce that \smash{$\big|A^{(\ell)}\big| \to 0$} uniformly.

By \eqref{eq:almostflat}, \smash{$\nabla^{S \otimes E^{(\ell)}} s^{(\ell)}$} is bounded in $L^2$. Using \eqref{eq:A_l_term}, we conclude that \smash{$\nabla^{S \otimes E} \big(\big(\id \otimes \sigma^{(\ell)}\big) s^{(\ell)}\big)$} is bounded in $L^2$. So, the sequence \smash{$\big(\id \otimes \sigma^{(\ell)}\big) s^{(\ell)} \in C^\infty(M,S \otimes E)$} is bounded in $H^1(M,S \otimes E)$.
After passing to a subsequence, the sequence \smash{$\big(\id \otimes \sigma^{(\ell)}\big) s^{(\ell)} \in C^\infty(M,S \otimes E)$} converges, in the weak topology of $H^1(M,S \otimes E)$, to a section $s$.
Since weak $H^1$-convergence implies strong $L^2$-convergence, the limit $s \in H^1(M,S \otimes E)$ satisfies
\[
\int_M \frac{1}{\rho(\Theta)} |s|^2 = 1.
\]
The inequality \eqref{eq:almostflat} implies that
\begin{equation}
\nabla_X^{S \otimes E} s + \frac{\rm i}{2} \Psi X \cdot s = 0
\label{eq:parallel}
\end{equation}
for every smooth vector field $X$ on $M$, where \eqref{eq:parallel} is understood in the sense of distributions.
Since \eqref{eq:parallel} holds for every smooth vector field $X$ on $M$, it follows
that $s$ is a weak solution of an overdetermined elliptic system.
By elliptic regularity, $s$ is smooth and \eqref{eq:parallel} holds classically.
Rescaling $s$ concludes the proof.
\end{proof}

\begin{Definition}
Let $t_0 \in S^1$, $f\colon M \to S^n$, and $E$ be defined as in Corollary~\ref{limiting_spinor_even_dim}.
We define the modified connection $\nabla^\Psi$ on $S \otimes E $ by
\begin{equation}
\nabla_X^\Psi s = \nabla_X^{S \otimes E} s + \frac{\rm i}{2} \Psi X \cdot s,
\label{def:nablapsi}
\end{equation}
where $\nabla^{S \otimes E}$ denotes the connection on $S \otimes E$ induced by the ones on $S$ and $E$.
\end{Definition}

\begin{Lemma}
The curvature tensor $R^\Psi$ of the connection $\nabla^\Psi$ defined in \eqref{def:nablapsi} satisfies:
\begin{align}
R^\Psi_{X,Y}s
&=
R^{S\otimes E}_{X,Y}s + \frac{\rm i}{2} ({\rm d}\Psi(X)Y-{\rm d}\Psi(Y)X) \cdot s - \frac{1}{4} \Psi^2 (X\cdot Y-Y\cdot X)\cdot s .
\label{eq:RPsi}
\end{align}
Here $R^{S\otimes E}$ denotes the curvature tensor of $\nabla^{S\otimes E}$.
Moreover, the curvature term in the Weitzen\-b\"ock formula satisfies
\begin{align}
\frac{1}{2} \sum_{\substack{1\leq j,k\leq n\\ j \neq k}} e_j\cdot e_k\cdot R^\Psi_{e_j,e_k}s
=
\left(\frac{1}{4} R + \mathcal{R}^E\right)s - \frac{{\rm i}(n-1)}{2} \nabla\Psi\cdot s + \frac{n(n-1)}{4} \Psi^2 s,
\label{eq:Weitzenbock}
\end{align}
where $e_1,\dots,e_n$ is a local orthonormal frame.
\end{Lemma}

\begin{proof}
We check \eqref{eq:RPsi} at a fixed point on $M$. Let $X$ and $Y$ be vector fields defined in a~neighborhood of that point whose covariant derivatives vanish at the point. We compute at that point:
\begin{align*}
\nabla^\Psi_X \nabla^\Psi_Y s
=
\nabla^{S \otimes E}_X \nabla^{S \otimes E}_Y s + \frac{\rm i}{2} {\rm d}\Psi(X)Y\!\cdot\! s + \frac{\rm i}{2} \Psi Y \!\cdot\! \nabla^{S \otimes E}_X s + \frac{\rm i}{2} \Psi X \!\cdot\! \nabla^{S \otimes E}_Y s
 - \frac{1}{4} \Psi^2 X \!\cdot\! Y\cdot s.
\end{align*}
Anti-symmetrizing with respect to $X$ and $Y$ yields \eqref{eq:RPsi}.

As to \eqref{eq:Weitzenbock}, we use formula~(8.8) in \cite[Chapter~II]{Lawson-Michelsohn} and \eqref{eq:RPsi} and we find
\begin{align*}
&\left(\frac{1}{4} R + \mathcal{R}^E\right)s =
\frac{1}{2} \sum_{j \neq k} e_j\cdot e_k\cdot R^{S \otimes E}_{e_j,e_k} s \\
&\qquad=
\frac{1}{2} \sum_{j \neq k} e_j\cdot e_k\cdot \left( R^\Psi_{e_j,e_k} s - \frac{\rm i}{2} ({\rm d}\Psi(e_j)e_k-{\rm d}\Psi(e_k)e_j) \cdot s + \frac{1}{4} \Psi^2 (e_j\cdot e_k-e_k\cdot e_j) \cdot s \right) \\
&\qquad=
\frac{1}{2} \sum_{j \neq k} e_j\cdot e_k\cdot R^\Psi_{e_j,e_k} s + \frac{{\rm i}(n-1)}{2} \nabla\Psi \cdot s - \frac{n(n-1)}{4} \Psi^2 s.\tag*{\qed}
\end{align*} \renewcommand{\qed}{}
\end{proof}

After these preparations, we now complete the proof of Theorem~\ref{spherical_band} for even $n$.
%\begin{proof}[Proof of Theorem~\ref{spherical_band} for $n$ even]
Let $t_0 \in S^1$, $f\colon M \to S^n$, and $E$ be defined as in Corollary~\ref{limiting_spinor_even_dim} and let $\nabla^\Psi$ denote the connection defined in~\eqref{def:nablapsi}.
By Corollary~\ref{limiting_spinor_even_dim}, there exists a section $s \in C^\infty(M,S \otimes E)$ such that $\int_M |s|^2 = 1$ and~${\nabla^\Psi s = 0}$.
Since $M$ is connected and $s$ is $\nabla^\Psi$-parallel, we have that $s \neq 0$ at each point in~$M$. Using \eqref{eq:Weitzenbock} and the fact that $R^\Psi$ annihilates $s$, we find
\begin{align*}
\mathcal{R}^E s
&=
-\frac{1}{4} R s + \frac{{\rm i}(n-1)}{2} \nabla \Psi \cdot s - \frac{n(n-1)}{4} \Psi^2 s \\
&=
-\frac{1}{4} R s + \frac{{\rm i}(n-1)}{2} \psi'(\Theta) \nabla\Theta \cdot s - \frac{n(n-1)}{4} \psi^2(\Theta) s .
\end{align*}
By Lemma~\ref{gradient_of_Theta}, $|\nabla \Theta| \leq 1$. Since $-\psi' > 0$ on $[\theta_-,\theta_+]$, it follows that{\samepage
\begin{align}
\big\langle \mathcal{R}^E s,s \big\rangle
&\leq
-\frac{1}{4} R |s|^2 - \frac{n-1}{2} \psi'(\Theta) |s|^2 - \frac{n(n-1)}{4} \psi^2(\Theta) |s|^2\nonumber \\
&\leq
-\frac{(n-2)(n-1)}{4} \frac{1}{\rho^2(\Theta)} |s|^2. \label{estimate_curvature_E}
\end{align}
In the last step, we have again used \eqref{ineq_for_scal_even_dim}.}

Let us fix an arbitrary point $x \in M$. We consider the singular values $\mu_1, \ldots, \mu_n \geq 0$ of ${\rm d}f_x\colon (T_x M,g) \to (T_{f(x)} S^n,g_{S^n})$, arranged in decreasing order. Since $f$ factors through $S^{n-1}$, the differential ${\rm d}f_x$ has rank at most $n-1$, and we obtain $\mu_n = 0$. The eigenvalues of the symmetric bilinear form $f^* g_{S^n}$ with respect to the metric $g$ are given by $\mu_1^2,\hdots,\mu_{n-1}^2,0$.

Using \eqref{ineq_for_g} together with the inequality $h(\cdot,t_0)^* g_{S^n} \leq g_{S^{n-1}}$, we obtain
\[g \geq \rho^2(\Theta) \varphi^* g_{S^{n-1}} \geq \rho^2(\Theta) f^* g_{S^n}\]
at the point $x$.
Therefore, $\mu_1^2,\hdots,\mu_{n-1}^2 \leq \frac{1}{\rho^2(\Theta)}$.
Using Proposition~\ref{curvature.term} in Appendix~\ref{appendix:Weitzen}, we deduce that
\[
 \big\langle \mathcal{R}^E s ,s \big\rangle \geq -\frac{(n-2)(n-1)}{4} \frac{1}{\rho^2(\Theta)} |s|^2
\]
at the point $x$.
Having established the reverse inequality in \eqref{estimate_curvature_E}, this inequality must be an equality.
In particular all the inequalities in \eqref{estimate_curvature_E} are equalities.
Since $n>2$ and $s \neq 0$ at the point $x$ and $\psi'\neq0$, we can draw the following conclusion:
\begin{itemize}\itemsep=0pt
\item The singular values of ${\rm d}f_x\colon (T_x M,g) \to (T_{f(x)} S^n,g_{S^n})$ are given by $\frac{1}{\rho(\Theta)}, \ldots, \frac{1}{\rho(\Theta)}, 0$. In other words, the eigenvalues of $f^* g_{S^n}$ with respect to $g$ at the point $x$ are given by $\frac{1}{\rho^2(\Theta)}, \ldots, \frac{1}{\rho^2(\Theta)}, 0$.
\item $|\nabla \Theta| = 1$ at the point $x$.
\end{itemize}
Since $|\nabla \Theta| = 1$ at the point $x$, Lemma~\ref{gradient_of_Theta} implies that ${\rm d}\varphi_x(\nabla \Theta) = 0$, hence ${\rm d}f_x(\nabla \Theta) = 0$.
Consequently, $\nabla \Theta$ lies in the nullspace of $f^* g_{S^n}$. Putting these facts together, we obtain
\[f^* g_{S^n} = \frac{1}{\rho^2(\Theta)} (g - {\rm d}\Theta \otimes{\rm d}\Theta),\]
hence
\begin{equation}
\label{a}
g = {\rm d}\Theta \otimes {\rm d}\Theta + \rho^2(\Theta) f^* g_{S^n}
\end{equation}
at the point $x$. On the other hand, using \eqref{ineq_for_g} together with the inequality $h(\cdot,t_0)^* g_{S^n} \leq g_{S^{n-1}}$, we obtain
\begin{equation}
\label{b}
g \geq {\rm d}\Theta \otimes {\rm d}\Theta + \rho^2(\Theta) \varphi^* g_{S^{n-1}} \geq {\rm d}\Theta \otimes {\rm d}\Theta + \rho^2(\Theta) f^* g_{S^n}
\end{equation}
at the point $x$. Combining (\ref{a}) and (\ref{b}), we conclude that
\[g = {\rm d}\Theta \otimes {\rm d}\Theta + \rho^2(\Theta) \varphi^* g_{S^{n-1}}\]
at the point $x$. Since $x$ is arbitrary, we conclude that $g= \Phi^*(g_0)$.
This means that $\Phi$ is a local isometry.
Since the target of $\Phi$ is simply connected and the domain is connected, it follows that~$\Phi$ is a global Riemannian isometry.
The proof of Theorem~\ref{spherical_band} for even $n$ is complete.

\subsection[Proof of Theorem~\ref{spherical_band} for n odd]{Proof of Theorem~\ref{spherical_band} for $\boldsymbol{n}$ odd}

When $n$ is odd, the proof of Theorem~\ref{spherical_band} is simpler, and we just indicate the necessary changes.
Instead of working with $\tilde{M} = M \times S^1$, we work with $M$.
Furthermore, instead of the map $\tilde{f} = h \circ (\varphi \times \id)\colon \tilde{M} \to S^n$, we directly work with $\varphi\colon M \to S^{n-1}$.
Let $S$ denote the spinor bundle over $M$, and let $E_0$ denote the spinor bundle over the round sphere $S^{n-1}$.
Since $n-1$ is even, we may decompose $E_0 = E_0^+ \oplus E_0^-$, where $E_0^+$ and $E_0^-$ denote the $\pm 1$-eigenbundles of the complex volume form.

\begin{Proposition}[cf.\ Cecchini--Zeidler \cite{Cecchini-Zeidler}]
\label{index_odd_dim}
Consider the indices of the following operators:
\begin{itemize}\itemsep=0pt
\item Let $\mathrm{ind}_1$ denote the index of the Dirac operator on $S \otimes \varphi^* E_0^+$ with boundary conditions $s = -{\rm i} \nu \cdot s$ on $\partial_+ M$ and $s = {\rm i} \nu \cdot s$ on $\partial_- M$.
\item Let $\mathrm{ind}_2$ denote the index of the Dirac operator on $S \otimes \varphi^* E_0^+$ with boundary conditions $s = {\rm i} \nu \cdot s$ on $\partial_+ M$ and $s = -{\rm i} \nu \cdot s$ on $\partial_- M$.
\item Let $\mathrm{ind}_3$ denote the index of the Dirac operator on $S \otimes \varphi^* E_0^-$ with boundary conditions $s = -{\rm i} \nu \cdot s$ on $\partial_+ M$ and $s = {\rm i} \nu \cdot s$ on $\partial_- M$.
\item Let $\mathrm{ind}_4$ denote the index of the Dirac operator on $S \otimes \varphi^* E_0^-$ with boundary conditions $s = {\rm i} \nu \cdot s$ on $\partial_+ M$ and $s = -{\rm i} \nu \cdot s$ on $\partial_- M$.
\end{itemize}
Then $\mathrm{ind}_1+\mathrm{ind}_2=0$, $\mathrm{ind}_3+\mathrm{ind}_4=0$, and $\max \{\mathrm{ind}_1,\mathrm{ind}_2, \mathrm{ind}_3, \mathrm{ind}_4\} > 0$.
\end{Proposition}

The proof of Proposition \ref{index_odd_dim} is analogous to the proof of Proposition \ref{index_even_dim} and uses the holographic index theorem.

After switching the bundles $E_0^+$ and $E_0^-$ if necessary, we may assume that ${\max \{\mathrm{ind}_1,\mathrm{ind}_2\} \!> 0}$.
As above, we focus on the case $\mathrm{ind}_1 > 0$.
(The case $\mathrm{ind}_2 > 0$ can be handled analogously.)

Let $E$ denote the pull-back of $E_0^+$ under $\varphi$. Let $\mathcal{D}^{S \otimes E}$ denote the Dirac operator on sections of $S \otimes E$, and let $\mathcal{D}^{\partial M}$ denote the boundary Dirac operator.

Let $\Psi$ be a smooth function on $M$ that will be specified later. If $s$ is section of the bundle~${S \otimes E}$ and $X$ is a vector field on $M$, we define a perturbed covariant derivative of $s$ by the formula
\[
 P_X s = \nabla_{X}^{S \otimes E}s + \frac{\rm i}{2} \Psi X \cdot s.
\]

\begin{Proposition}
\label{integral_formula_odd_dim}
Let $s \in C^\infty(M,S \otimes E)$.
Then
\begin{align*}
&- \int_M \left|\mathcal{D}^{S \otimes E} s - \frac{{\rm i}n}{2} \Psi s\right|^2 + \int_M |Ps|^2 + \frac{1}{4} \int_M R |s|^2 \\
&+ \int_M \big\langle \mathcal{R}^E s,s \big\rangle + \frac{n(n-1)}{4} \int_M \Psi^2 |s|^2 - \frac{{\rm i}(n-1)}{2} \int_M \langle (\nabla \Psi) \cdot s,s \rangle \\
&\qquad= \frac{1}{2} \int_{\partial_+ M} \big\langle \mathcal{D}^{\partial M} s,s + {\rm i} \nu \cdot s \big\rangle + \frac{1}{2} \int_{\partial_+ M} \big\langle s + {\rm i} \nu \cdot s,\mathcal{D}^{\partial M} s \big\rangle \\
&\phantom{\qquad=}{} - \frac{1}{2} \int_{\partial_+ M} (H-(n-1)\Psi) |s|^2 - \frac{n-1}{2} \int_{\partial_+ M} \Psi \langle s + {\rm i} \nu \cdot s,s \rangle \\
&\phantom{\qquad=}{} + \frac{1}{2} \int_{\partial_- M} \big\langle \mathcal{D}^{\partial M} s,s - {\rm i} \nu \cdot s \big\rangle + \frac{1}{2} \int_{\partial_- M} \big\langle s - {\rm i} \nu \cdot s,\mathcal{D}^{\partial M} s \big\rangle \\
&\phantom{\qquad=}{} - \frac{1}{2} \int_{\partial_- M} (H+(n-1)\Psi) |s|^2 + \frac{n-1}{2} \int_{\partial_- M} \Psi \langle s - {\rm i} \nu \cdot s,s \rangle.
\end{align*}
\end{Proposition}

The proof of Proposition~\ref{integral_formula_odd_dim} is analogous to Proposition~\ref{integral_formula_even_dim}.

As above, we define a function $\psi\colon [\theta_-,\theta_+] \to \R$ by $\psi(\theta) = \frac{\rho'(\theta)}{\rho(\theta)}$. The assumption $R \geq R_{g_0} \circ \Phi$ gives
\begin{equation} \label{ineq_for_scal_odd_dim}
R \geq (n-1) \left( -2 \psi'(\Theta) - n \psi^2(\Theta) + \frac{n-2}{\rho^2(\Theta)} \right).
\end{equation}
We define $\Psi\colon M \to \R$ by $\Psi = \psi \circ \Theta$.

At this point, we use the assumption that $\mathrm{ind}_1 > 0$. In view of the deformation invariance of the index, we find a section $s$ of $S \otimes E$ such that
\begin{itemize}\itemsep=0pt
 \item $s$ does not vanish identically,
 \item $D^{S \otimes E } s - \frac{{\rm i}n}{2} \psi s = 0$ on $M$,
 \item $s = -{\rm i} \nu \cdot s$ on $\partial_+ M$ and $s = {\rm i} \nu \cdot s$ on $\partial_- M$.
\end{itemize}
Using Proposition~\ref{integral_formula_odd_dim}, we obtain
\begin{align}
&\int_M |P s|^2 + \frac{1}{4} \int_M R |s|^2
+ \int_M \big\langle \mathcal{R}^E s,s \big\rangle + \frac{n(n-1)}{4} \int_M \Psi^2 |s|^2 - \frac{{\rm i}(n-1)}{2} \int_M \langle (\nabla \Psi) \cdot s,s \rangle\nonumber \\
&\qquad= -\frac{1}{2} \int_{\partial_+ M} (H-(n-1)\Psi) |s|^2 - \frac{1}{2} \int_{\partial_- M} (H+(n-1)\Psi) |s|^2. \label{important_estimate_odd_dim}
\end{align}
By assumption, $H-(n-1)\Psi \geq 0$ on $\partial_+ M$ and $H +(n-1)\Psi \geq 0$ on $\partial_- M$. Therefore, the right-hand side in \eqref{important_estimate_odd_dim} is non-positive.

Fix a point $x \in M $, and let $ \mu_1,\hdots,\mu_n \geq 0$ denote the singular values of the differential ${\rm d}\varphi_x\colon (T_x M,g ) \to \big(T_{f(x)} S^{n-1} ,g_{S^{n-1}}\big)$, arranged in decreasing order. Since the differential ${\rm d}\varphi_x$ has rank at most $n-1$, it follows that $\mu_n = 0$.

Since $\varphi$ is $1$-Lipschitz by assumption, we obtain
\[
 g \geq \rho^2(\Theta) \varphi^* g_{S^{n-1}} .
 \]
Consequently, $\mu_1^2,\hdots,\mu_{n-1}^2 \leq \frac{1}{\rho^2(\Theta)}$. In view of Proposition~\ref{curvature.term}, this implies
\begin{align*}
\big\langle \mathcal{R}^{E} s,s \big\rangle
&\geq -\frac{1}{4} \sum_{\substack{1 \leq j,k \leq n\\ j \neq k}} \mu_j \mu_k |s|^2 \geq -\frac{(n-2)(n-1)}{4} \frac{1}{\rho^2(\Theta)} |s|^2.
\end{align*}
Since $|\nabla \Theta| \leq 1$ and $|\nabla \Psi| \leq -\psi'(\Theta)$, this gives the pointwise estimate
\begin{align*}
&\frac{1}{4} R |s|^2 + \big\langle \mathcal{R}^{E} s,s \big\rangle + \frac{n(n-1)}{4} \Psi^2 |s|^2 - \frac{n-1}{2} |\nabla \Psi| |s|^2 \\
&\qquad\geq \frac{1}{4} R |s|^2 - \frac{(n-2)(n-1)}{4} \frac{1}{\rho^2(\Theta)} |s|^2
+ \frac{n(n-1)}{4} \psi^2(\Theta) |s|^2 + \frac{n-1}{2} \psi'(\Theta) |s|^2 \geq 0.
\end{align*}
In the last step, we have used \eqref{ineq_for_scal_odd_dim}.
Putting these facts together, we conclude that
\[
\int_M | P s|^2 = 0,
\]
hence
\[
 \nabla_X^{S \otimes E} s + \frac{\rm i}{2} \psi(\Theta) X \cdot s = 0
\]
for every vector field $X$.
This is the analogue of Corollary~\ref{limiting_spinor_even_dim}.
From here on, the proof of Theorem~\ref{spherical_band} in the odd-dimensional case proceeds in the same way as in the even-dimensional case.
